\documentclass[12pt]{article}
\usepackage[T1]{fontenc}
\usepackage[utf8]{inputenc}
\usepackage{lmodern}
\usepackage{textcomp}
\usepackage{enumitem}
\usepackage{geometry}
\usepackage{graphicx}
\usepackage{amsmath, amssymb}
\geometry{margin=1in}
\begin{document}
\begin{center}
Economics - Graduate Level Assessment 8 \
Total Marks: 40 \
Answer all questions.
\end{center}

\section*{Multiple Choice (10 marks)}
\begin{enumerate}[label=\arabic*.]
\item Why is the danger inherent in the issuance of fiat money?
\begin{enumerate}[label=(\alph*)]
    \item It is often used to meet emergencies.
    \item Its quantity is controlled by the government rather than by market forces.
    \item It is backed by a physical commodity like gold or silver.
    \item It encourages the accumulation of debt through interest-free loans.
    \item It is used for buying military supplies during a war.
\end{enumerate}
\item Production possibilities frontiers are concave to the origin because
\begin{enumerate}[label=(\alph*)]
    \item of the principle of comparative advantage.
    \item of the existence of fixed resources.
    \item of the scarcity of resources.
    \item of the law of decreasing costs.
    \item of inefficiencies in the economy.
\end{enumerate}
\item At a Nash equilibrium,
\begin{enumerate}[label=(\alph*)]
    \item the equilibrium is unstable and each party would like to switch strategies
    \item the marginal utility curve intersects the marginal cost curve
    \item neither party has an incentive to deviate from his or her strategy
    \item the equilibrium is stable and each party would like to switch strategies
    \item the equilibrium is where the total revenue equals total cost
\end{enumerate}
\item What might happen to the housing construction industry if com-mercial banks are allowed to compete freely with mutual sav-ings banks?
\begin{enumerate}[label=(\alph*)]
    \item The quality of housing construction would deteriorate
    \item More houses could be bought
    \item The housing industry would slump
    \item Mutual savings banks would become obsolete
    \item Deposits in mutual savings banks would increase
\end{enumerate}
\item A negative or contractionary supply shock will
\begin{enumerate}[label=(\alph*)]
    \item shift the Phillips curve to the left.
    \item shift the investment demand curve to the right.
    \item shift the money demand curve to the right.
    \item shift the Phillips curve to the right.
\end{enumerate}
\end{enumerate}

\section*{True / False (10 marks)}
\textit{Answer True or False}
\begin{enumerate}[label=\arabic*.]
\setcounter{enumi}{5}
\item True or False: In monopolistic competition, firms are price takers and cannot influence the market price of their products.
\item True or False: The demand curve for a perfectly competitive firm's product is perfectly elastic, reflecting the price-taking nature of such firms.
\item True or False: The optimal one-step ahead forecast of the change in y is the average value of the change in y over the in-sample period.
\item True or False: Continuously compounded returns can be interpreted as continuously compounded changes in the prices of assets.
\item True or False: A monopolist sets price equal to average total cost, leading to zero economic profit in the long run.
\end{enumerate}

\section*{Long Form Response (20 marks)}
\textit{Answer each question with a clear and complete explanation.}
\begin{enumerate}[label=\arabic*.]
\setcounter{enumi}{10}
\item Discuss the implications of a capital/output ratio of 2.5 on investment levels when GNP increases from $985 billion to $1,055 billion, and examine the effect of a marginal propensity to consume of 0.5 on GNP changes.
\item Analyze a scenario where a commercial bank has $100,000 in deposits and $65,000 in loans with a required reserve ratio of 20\%. Discuss how much additional money the bank can lend and evaluate the potential increase in the overall money supply resulting from these loans.
\end{enumerate}

\vfill
\noindent\textit{End of Paper}
\end{document}